% ======================================================================
% PLANTILLA LaTeX — Euclides
% ----------------------------------------------------------------------
% Motor requerido: LuaLaTeX o XeLaTeX (se usa fontspec para la fuente).
%   Compilar con:  lualatex nombre.tex   (dos veces, por las referencias
%                  cruzadas de ecuaciones, definiciones, teoremas, etc.)
% ======================================================================
\documentclass[11pt,a4paper,twoside,openright]{book}    
\usepackage{amsmath,amsfonts,amssymb,amsthm,amscd} 


% se copian las variables más relevantes
\makeatletter
\let\newtitle\@title
\let\newauthor\@author
\let\newdate\@date
\makeatother


% ----------------------------------------------------------------------
% Fuente y tamaño de letra
% ----------------------------------------------------------------------
\usepackage{fontspec}
\usepackage{newcomputermodern} % Tipo de letra, New Computer Modern

%
%% Tipo de letra de las letras capitales.
%

%\newfontfamily\chapnumfont{arenq.regular}[Path = ./, Extension = .otf] % NO
%\newfontfamily\chapnumfont{manbow.lines-regular}[Path = ./, Extension = .otf] % de momento, la que más me gusta
%\newfontfamily\chapnumfont{evensong}[Path = ./, Extension = .otf]
%\newfontfamily\chapnumfont{berylium.bold}[Path = ./, Extension = .otf]
%\newfontfamily\chapnumfont{library-3-am.3am}[Path = ./, Extension = .otf] % este 
%\newfontfamily\chapnumfont{foglihtenno04.regular}[Path = ./, Extension = .otf]
%\newfontfamily\chapnumfont{labyrinth.regular}[Path = ./, Extension = .otf]
%\newfontfamily\chapnumfont{baflion-sans.black}[Path = ./, Extension = .otf]
%\newfontfamily\chapnumfont{atavian.regular}[Path = ./, Extension = .otf]
%\newfontfamily\chapnumfont{stripescaps.regular}[Path = ./, Extension = .ttf] % non me gusta e non ten tilde
%\newfontfamily\chapnumfont{dita-sweet.regular}[Path = ./, Extension = .otf]
%\newfontfamily\chapnumfont{ZnikomitNo25}[Path = ./, Extension = .otf]
\newfontfamily\chapnumfont{Archistico_Bold}[Path = ./, Extension = .ttf]
%\newfontfamily\chapnumfont{Roke1984}[Path = ./, Extension = .otf]
%\newfontfamily\chapnumfont{SketchiquaB}[Path = ./, Extension = .ttf]
%\newfontfamily\chapnumfont{WelcomeToFuture1-E4Xnl}[Path = ./, Extension = .ttf]
%\newfontfamily\chapnumfont{CrazyConstructor}[Path = ./, Extension = .ttf]
%\newfontfamily\chapnumfont{Limelight-Regular}[Path = ./, Extension = .ttf]



% interlineado
\renewcommand{\baselinestretch}{1.25}
\setlength{\parskip}{1mm}
\raggedbottom\setlength{\baselineskip}{1mm}
\allowdisplaybreaks



% ----------------------------------------------------------------------
% Idioma
% ----------------------------------------------------------------------
\usepackage[spanish,es-lcroman,es-noquoting]{babel} % Español



% ----------------------------------------------------------------------
% Color primario
% ----------------------------------------------------------------------
\usepackage{xcolor}
\definecolor{colorprimario}{HTML}{FF4000}%{FF461A}



% ----------------------------------------------------------------------
% Geometría de la Página
%-----------------------------------------------------------------------
\usepackage[papersize={210mm,297mm}, left=2.5cm, right=2.5cm, top=3.5cm, bottom=3.5cm]{geometry}


% ----------------------------------------------------------------------
% Paquetes auxiliares
% ----------------------------------------------------------------------
\usepackage{array}       % para el separador "!{...}" en tablas
\usepackage{tabularx}     % columna elástica "X" 
\usepackage{framed}       % entorno "leftbar": regla vertical izquierda
\usepackage{graphicx}
\usepackage{caption}
\usepackage{tikz}
\usepackage{calc}
\usepackage{ulem} % para poder subrayar correctamente
\usepackage{scrextend} % para comandos como addmargin
\usepackage{needspace} % para evitar roturas de cajas no bonitas



% ----------------------------------------------------------------------
% Apéndices
% ----------------------------------------------------------------------
%% Formato: 
%% - Los apéndices se disponen en una Parte separada, 
%%
%% - con una página indicando el inicio con `\appendixtitle` 
%%   escrito con la primera letra en mayúscula y de `colorprimario`. 
%%
%% - El título está subrayado y centrado en el
%%   centro de la página de estilo `empty`, en una página a la izquierda.
%%
%% - Siguiente página en blanco.
%

\usepackage[toc,page]{appendix}

%
%% comandos de formateo
%

\newcommand*{\appendixfontsize}{50pt} % tamaño del título de la sección de anexos
\newcommand*{\appendixtitle}{Apéndices} % título de la sección de Anexos


% TRABAJO INTERNO: -----------------------------------------------------
\newlength{\appendixfontsizeInternal}

%
%% Formato del título de la zona de apéndices
%
\makeatletter
% 1. El motor interno: toma la primera letra (#1), el resto (#2), y se detiene en \@nil
\def\@appendixformat#1#2\@nil{%
	% Primera letra (#1) con su estilo y subrayado propio
	\uline{{\fontsize{1.8\appendixfontsizeInternal}{1.8\appendixfontsizeInternal}\chapnumfont\selectfont\bfseries\color{colorprimario}#1}%
	% Resto del texto (#2) con su propio estilo
	{\fontsize{\appendixfontsize}{\appendixfontsize}\selectfont\bfseries#2}}%
}

% 2. El comando público: expande tu variable antes de enviarla al motor interno
\newcommand{\formattingappendixname}[1]{%
	\expandafter\@appendixformat#1\@nil
}

% 
%% Definición del apéndice como tal
%
\gdef\appendix{%
	\setlength{\appendixfontsizeInternal}{\appendixfontsize}
	\renewcommand{\ULthickness}{1.5pt}
	\cleardoublepage
	\thispagestyle{empty} % Quita encabezados y pies de página en esta portada
	
	% Página con "Apéndices" alineado a la izquierda y centrado verticalmente
	\vspace*{\fill}
	\begin{center}
		\formattingappendixname{\appendixtitle} % formateo del título
	\end{center}
	\vspace*{\fill}
	
	\clearpage
	\thispagestyle{empty}
	\null % Inserta una página en blanco obligatoria
	\clearpage
	
	% Restablece los contadores para los capítulos del apéndice
	\setcounter{chapter}{0}%
	\setcounter{section}{0}%
	\gdef\@chapapp{\appendixname}%
	\gdef\thechapter{\@Alph\c@chapter}%
}
\makeatletter



% ----------------------------------------------------------------------
% Índice de contenidos
% ----------------------------------------------------------------------
\usepackage{titletoc}

%
%% comandos de formato
%

\renewcommand\contentsname{Índice} % nombre del título

\newcommand{\contentstitlesize}{50pt} % tamaño de letra del título




% TRABAJO INTERNO ------------------------------------------------------
\newlength{\contentstitlesizeInternal}


\titlecontents{chapter}[0pt]
{\addvspace{1em}\bfseries}                 % antes de la entrada
{\textcolor{colorprimario}\thecontentslabel\quad}        % con número
{}                                            % sin número (no se usa aquí)
{\leavevmode                                  % <- después del título:
	\hfill\kern0pt    %    
	\textcolor{colorprimario}{\thecontentspage}} %    página en naranja


\titlecontents{section}[1.5em]
{\addvspace{0.10em}}
{\textcolor{colorprimario}\thecontentslabel.\quad}
{}
{\leavevmode
	~\leaders\hrule height 0.5pt\hfill\kern0pt ~   % línea recta de relleno
	\textcolor{colorprimario}{\thecontentspage}}


\titlecontents{subsection}[3em]
{\addvspace{0.1em}\small}
{\thecontentslabel.\quad}
{}
{\leavevmode
	~\leaders\hrule height 0.5pt\hfill\kern0pt ~
	\textcolor{colorprimario}{\thecontentspage}}




\let\oldtableofcontents\tableofcontents % guardamos la anterior forma de table of contents.

% Formateo del título
\makeatletter
\renewcommand{\tableofcontents}{%
	\begingroup
	\hypersetup{linkcolor=black}% cambiar, de forma temporal, el color de los títulos.
	\renewcommand{\@makeschapterhead}[1]{%
		\setlength{\contentstitlesizeInternal}{\contentstitlesize}
		\renewcommand{\arrayrulewidth}{1.5pt}
		\vspace*{-5mm}%
		\noindent
		\begin{tabularx}{\textwidth}{ l @{\vline \phantom{M}} X}
			 & \vphantom{{\fontsize{1.8\contentstitlesizeInternal}{1.8\contentstitlesizeInternal}\selectfont M}} \textbf{{\fontsize{1.2\contentstitlesizeInternal}{1.2\contentstitlesizeInternal}\chapnumfont\selectfont\color{colorprimario}Í}\scshape\fontsize{\contentstitlesizeInternal}{\contentstitlesizeInternal}\selectfont ndice} \\
			\llap{\raisebox{1ex}{\rule{20cm}{\arrayrulewidth}}\hspace{-\textwidth}} &\\ &
		\end{tabularx}
		\renewcommand{\arrayrulewidth}{1pt}
	}%
	\oldtableofcontents
	\endgroup
}
\makeatother



% ----------------------------------------------------------------------
% Cabecera: 
% ----------------------------------------------------------------------
\usepackage{fancyhdr}
\pagestyle{fancy}
\fancyhf{}
\renewcommand{\headrulewidth}{0.5pt}
\renewcommand{\footrulewidth}{0pt}

\renewcommand{\chaptermark}[1]{\markboth{\textbf{{\color{colorprimario}\thechapter} #1}}{}} % formato del título de capítulo

\fancypagestyle{fancy}{	
	\fancyhf{}
	\renewcommand{\headrulewidth}{0.5pt}
	\renewcommand{\footrulewidth}{0pt}
	\fancyhead[RO]{\bfseries\fontsize{11pt}{13pt}\selectfont \newtitle}
	\fancyhead[LE]{\fontsize{11pt}{13pt}\selectfont \leftmark}
	\fancyhead[RE,LO]{\color{colorprimario}\bfseries\thepage}
}

\fancypagestyle{plain}{%
	\fancyhf{}
	\renewcommand{\headrulewidth}{0pt}
%	\renewcommand{\footrulewidth}{0pt}
	\fancyfoot[C]{\color{colorprimario}\bfseries\thepage}
}


% ----------------------------------------------------------------------
% Títulos de capítulo y secciones
% ----------------------------------------------------------------------
\usepackage{titlesec}


% titulos de secciones
\titleformat{\section}
{\normalfont\Large\bfseries}
{\color{colorprimario}\thesection}
{0.5em}
{\hyphenpenalty=10000\exhyphenpenalty=10000\relax}

\titleformat{\subsection}
{\normalfont\large\bfseries}
{\color{colorprimario}\thesection.\arabic{subsection}}
{0.5em}
{\hyphenpenalty=10000\exhyphenpenalty=10000\relax}


% títulos de capítulo
\makeatletter

%
%% Comandos de Formato
%

\newcommand*{\chaptertitlesize}{25pt} % definición del tamaño de letra en los títulos
\newcommand*{\chapternumsize}{90pt} % definición del tamaño de fuenta del número de capítulo

\newcommand{\chaptersubsep}{10mm}

% TRABAJO INTERNO ------------------------------------------------------
\newlength{\chaptertitlesizeInternal}
\setlength{\chaptertitlesizeInternal}{\chaptertitlesize}

\newlength{\chapternumsizeInternal}
\setlength{\chapternumsizeInternal}{\chapternumsize}

\newlength{\chaptersubsepInternal}
\setlength{\chaptersubsepInternal}{\chaptersubsep}

\renewcommand{\tabularxcolumn}[1]{b{#1}}

\renewcommand{\@makechapterhead}[1]{%
	\setlength{\chaptertitlesizeInternal}{\chaptertitlesize}
	\setlength{\chapternumsizeInternal}{\chapternumsize}
	\setlength{\chaptersubsepInternal}{\chaptersubsep}
	\renewcommand{\arrayrulewidth}{1.5pt}
	\noindent
	\begin{tabularx}{\textwidth}{ l @{\vline \phantom{M}} X}
%		\llap{\raisebox{1ex}{\rule{25cm}{\arrayrulewidth}}\hspace{-0pt-\widthof{\fontsize{90}{90}\chapnumfont\selectfont\thechapter}}\hspace{-3mm}} & \\ % línea superior en el número
		{\fontsize{\chapternumsize}{\chapternumsize}\chapnumfont\selectfont \color{colorprimario} \thechapter}
		 \vphantom{{\fontsize{1.3\chapternumsizeInternal}{1.3\chapternumsizeInternal}\selectfont \thechapter}} 
		 & \textbf{\scshape\fontsize{\chaptertitlesize}{\chaptertitlesize}\selectfont #1} \\
		\llap{\raisebox{1ex}{\rule{20cm}{\arrayrulewidth}}\hspace{-\textwidth}} & \\ &
	\end{tabularx}
	\vspace{\chaptersubsepInternal}
	\renewcommand{\arrayrulewidth}{1pt}
}

% título sin numerar
\renewcommand{\@makeschapterhead}[1]{%
	\setlength{\chaptertitlesizeInternal}{\chaptertitlesize}
	\setlength{\chapternumsizeInternal}{\chapternumsize}
	\setlength{\chaptersubsepInternal}{\chaptersubsep}
	\renewcommand{\arrayrulewidth}{1.5pt}
	\par\noindent
	\begin{tabularx}{\textwidth}{ l @{\vline \phantom{M}} X}
	%	{\fontsize{90}{90}\selectfont \color{colorprimario} \thechapter} 
	%\llap{\raisebox{1ex}{\rule{25cm}{\arrayrulewidth}}\hspace{-3mm}} & \\
		 \vphantom{{\fontsize{\dimexpr 25pt+\chaptertitlesizeInternal\relax}{\dimexpr 25pt+\chaptertitlesizeInternal\relax}\selectfont M}}
		  & \textbf{\scshape\fontsize{\chaptertitlesizeInternal}{\chaptertitlesizeInternal}\selectfont #1} \\
		\llap{\raisebox{1ex}{\rule{20cm}{\arrayrulewidth}}\hspace{-\textwidth}} &\\ &
	\end{tabularx}
	\vspace{\chaptersubsepInternal}
	\renewcommand{\arrayrulewidth}{1pt}
}
\makeatother


%\titlespacing*{\chapter}{0pt}{-1.5cm}{2cm}
\titlespacing*{\section}{0pt}{0.25cm}{0.25cm}



% ----------------------------------------------------------------------
% Notas al pie coloreadas
% ----------------------------------------------------------------------
\renewcommand{\thefootnote}{\textcolor{colorprimario}{\bfseries\arabic{footnote}}}


% ----------------------------------------------------------------------
% Enlaces coloreados y en negrita
% ----------------------------------------------------------------------
\usepackage[colorlinks=true, linkcolor=colorprimario, citecolor=colorprimario,
urlcolor=colorprimario, filecolor=colorprimario,backref=page]{hyperref}
\let\oldhref\href
\renewcommand{\href}[2]{\oldhref{#1}{\textbf{#2}}}



% ----------------------------------------------------------------------
% Figuras y tablas
% ----------------------------------------------------------------------
\DeclareCaptionLabelFormat{coloreado}{\textcolor{colorprimario}{\textbf{#1~#2}}}
\captionsetup{
	labelformat=coloreado, 
	labelsep=colon}



% ----------------------------------------------------------------------
% Ecuaciones: numerar solo las que están referenciadas + número coloreado
% ----------------------------------------------------------------------
% \usepackage{mathtools}


% IMPORTANTE: el tagform coloreado se define ANTES de activar
% showonlyrefs. mathtools guarda una copia de \tagform@ en el momento
% de activar showonlyrefs y la envuelve con su lógica de "solo
% mostrar si está referenciada"; si se redefine \tagform@ DESPUÉS, se
% destruye esa lógica en vez de envolverla (y el número sale mal).
\makeatletter
\renewcommand{\tagform@}[1]{\maketag@@@{\color{colorprimario}\bfseries(#1)}}
\makeatother



% ----------------------------------------------------------------------
% Bibliografía
% ----------------------------------------------------------------------
% Cambia el separador cuando solo hay dos páginas (ej: 1 y 2)
\renewcommand*{\backreftwosep}{ y }

% Cambia el separador final cuando hay tres o más páginas (ej: 1, 2 y 3)
\renewcommand*{\backreflastsep}{ y }

% Desactiva el formato por defecto
\renewcommand*{\backref}[1]{}

% Define el nuevo formato
%\renewcommand*{\backrefalt}[4]{%
%	\ifcase #1 %
%	% Caso 0: No se ha citado en el texto (puedes dejarlo vacío)
%	\or
%	% Caso 1: Citado en exactamente una página
%	\\ \footnotesize \raggedleft  \textit{#2.}%
%	\else
%	% Caso 2: Citado en múltiples páginas
%	\par \hfill \footnotesize  \texttt{pp. #2.}%
%%	\begin{flushright}
%%		\footnotesize {\textit{Citado en las páginas #2.}}%
%%	\end{flushright} 
%	\fi
%}

% -----
% Titulo: backreferences a las páginas donde las referencias bibliográficas aparecen. 
% Formato:  recuadro grisáceo con las páginas en `colorprimario`
% -----
\makeatletter
\renewcommand*{\backrefalt}[4]{%
	\ifnum#1=0\else
	\nobreak\hfill\break
	\vspace{-25pt}%
	\begin{flushright}
		\def\@lastpage{}%
		\@for\@page:=#2\do{%
			\ifx\@page\@lastpage\else
			\begingroup
			\hypersetup{linkcolor=colorprimario}%
			\setlength{\fboxsep}{2pt}%
			\colorlet{backrefbg}{gray!15}
			\colorlet{backrefborder}{gray!55}
			\fcolorbox{backrefborder}{backrefbg}
			{\footnotesize\bfseries\@page}%
			\endgroup
			\hspace{3pt}%
			\let\@lastpage\@page
			\fi
		}%
	\end{flushright}%
	\fi
}
\makeatother


% cuadro al rededor de las referencias bibliográficas
%\newcommand{\boxcite}[1]{%
%	\begingroup
%	\setlength{\fboxsep}{2.5pt}%
%	\setlength{\fboxrule}{0.8pt}%
%	\fcolorbox{black!20}{white}{%
%		\textcolor{black}{\footnotesize\bfseries#1}%
%	}%
%	\endgroup
%}
%
%\makeatletter
%\renewcommand{\@cite}[2]{%
%	\boxcite{[#1\if@tempswa , #2\fi]}%
%}
%\makeatother

\emergencystretch=5em




% ----------------------------------------------------------------------
% Entornos: Definición, Teorema, Proposición, Lema
% ----------------------------------------------------------------------


% Regla vertical izquierda reutilizable (definiciones/teoremas/.../demos)
\renewenvironment{leftbar}{%
	\def\FrameCommand{{\color{black}\vrule width 0.5pt}\hspace{3mm}}%
	\vspace{-9mm}
	\MakeFramed{\advance\hsize-\width \FrameRestore}%
	\vphantom{M}
	\vspace{1mm}\\
}{\endMakeFramed}



\newlength{\anchocaja}
\newcommand{\openEuclidOverlayleft}[2]{%
	\par\addvspace{1.1em}
	\renewcommand{\arrayrulewidth}{0.5pt}
	\renewcommand{\arraystretch}{0.5}
	\setlength{\anchocaja}{\widthof{MM #1 | #2 MM}}
	\ifdim\anchocaja>\textwidth
	\setlength{\anchocaja}{\textwidth}
	\fi
	\setbox0=\hbox{%
		\begin{tabularx}{\anchocaja}{r|X}
			\vphantom{\Large M}~ \color{colorprimario}\bfseries #1 & #2 \\[0ex]
			\hline
			&
		\end{tabularx}%
	}%
	\Needspace{\dimexpr\ht0+\dp0+3\baselineskip\relax}
	\hspace{-2em}\box0
	\par\nobreak
%	\vskip -9mm\nobreak
	\begin{leftbar}
	}
	\newcommand{\closeEuclidOverlayleft}{%
	\end{leftbar}
	\renewcommand{\arraystretch}{1}
	\addvspace{1.1em}
}



% para las notas

\newlength{\anchocajanota}
\newcommand{\openEuclidOverlayright}[2]{%
	\par\addvspace{1.1em}
	\renewcommand{\arrayrulewidth}{0.5pt}
	\renewcommand{\arraystretch}{0.5}
	\setlength{\anchocajanota}{\textwidth}
%	\hspace{-2em}
	\setbox0=\hbox{
		\begin{tabularx}{\anchocajanota}{X|l}
			\flushright \color{colorprimario}\bfseries #1 & #2 \\[0ex]
			\hline
			&
		\end{tabularx}
	}
	\Needspace{\dimexpr\ht0+\dp0+2\baselineskip\relax}
	\hspace{-2.2em}\box0
	\par\nobreak
%	\nobreak\hfill
	\begin{leftbar}
	}
	\newcommand{\closeEuclidOverlayright}{%
	\end{leftbar}
	\renewcommand{\arraystretch}{1}
	\addvspace{1.1em}
}




% definición de los entornos
\usepackage{etoolbox} % para \ifstrequal

\newcommand{\newenv}[4]{%
	% #1 = identificador, #2 = ámbito, #3 = título, #4 = left|right
	\ifstrequal{#2}{chapter}{%
		\newcounter{#1}[chapter]%
		\expandafter\renewcommand\csname the#1\endcsname{\thechapter.\arabic{#1}}%
	}{%
		\ifstrequal{#2}{section}{%
			\newcounter{#1}[section]%
			\expandafter\renewcommand\csname the#1\endcsname{\thechapter.\arabic{section}.\arabic{#1}}%
		}{%
			\ifstrequal{#2}{subsection}{%
				\newcounter{#1}[subsection]%
				\expandafter\renewcommand\csname the#1\endcsname{\thechapter.\arabic{section}.\arabic{subsection}.\arabic{#1}}%
			}{%
				\newcounter{#1}%
				\expandafter\renewcommand\csname the#1\endcsname{\arabic{#1}}%
	}}}%
	\newenvironment{#1}[1][]{%
		\refstepcounter{#1}%
		\csname openEuclidOverlay#4\endcsname{#3~\csname the#1\endcsname}{##1}%
	}{\csname closeEuclidOverlay#4\endcsname}%
}



% ----------------------------------------------------------------------
% Demostración
% ----------------------------------------------------------------------
\renewcommand{\proofname}{Demostración}
\renewenvironment{proof}[1][\proofname]{%
	\addvspace{1em}\par
	\Needspace{3\baselineskip}
	\pushQED{\qed}%
	\noindent{\color{colorprimario}\bfseries #1}%
	\begin{addmargin}{3mm}
	}{%
		\popQED
	\end{addmargin}
	\addvspace{1em}
}


% ----------------------------------------------------------------------
% Blóques de código
% ----------------------------------------------------------------------

%\usepackage{fancyvrb}
%\fvset{
%	fontsize=\small,
%	frame=single,
%	framesep=3mm,
%	rulecolor=\color{black},  % color del marco
%	numbers=left,                      % numeración de líneas
%	numbersep=5pt,
%	formatcom=\color{black},           % color del texto en sí
%}

\usepackage{listings}

\lstdefinestyle{miR}{
	language=R,
	basicstyle=\ttfamily\small,
	keywordstyle=\color{colorprimario}\bfseries,   
	commentstyle=\color{gray}\itshape,  % # comentarios
	stringstyle=\color{blue}, % "cadenas"
	numbers=left, numbersep=6pt, numberstyle=\tiny\color{gray},
	showstringspaces=false, breaklines=true, xleftmargin=-1.2em,
}


%\usepackage{tcolorbox}
%\tcbuselibrary{listings, breakable}
%
%\newtcblisting{codigo}{
%	breakable,
%	colback=white, 
%	colframe=black,
%	boxrule=0.6pt, 
%	listing only,
%	listing options={style=miR, basicstyle=\ttfamily\small, breaklines=true},
%%	left=3mm, right=3mm, top=1mm, bottom=1mm,
%}


\usepackage{tikz}
\usepackage[many]{tcolorbox}
\tcbuselibrary{listings,skins, breakable}

% Definimos el nuevo entorno para la caja
\newtcblisting{codigo}{%
	breakable,
	enhanced,
	frame hidden,       % Ocultamos el marco normal
	colback=white,      % Color de fondo de la caja
	top=1mm, bottom=1mm, left=4mm, right=4mm, % Márgenes internos
	listing only,
	listing options={style=miR, basicstyle=\ttfamily\small, breaklines=true},
	overlay={
		\def\ext{2mm}     % Define cuántos milímetros sobresalen las líneas
		
		% Línea superior
		\draw[thick, black] ([xshift=-\ext]frame.north west) -- ([xshift=\ext]frame.north east);
		% Línea inferior
		\draw[thick, black] ([xshift=-\ext]frame.south west) -- ([xshift=\ext]frame.south east);
		% Línea izquierda
		\draw[thick, black] ([yshift=\ext]frame.north west) -- ([yshift=-\ext]frame.south west);
		% Línea derecha
		\draw[thick, black] ([yshift=\ext]frame.north east) -- ([yshift=-\ext]frame.south east);
	},
}



% ----------------------------------------------------------------------
% glosario
% ----------------------------------------------------------------------
\usepackage[toc]{glossaries}
\setglossarystyle{listhypergroup}





